\subsection{Pregunta 2: verificación global de conectividad}
\autor{Bárbara}

\textbf{Pruebas pedidas:} ping y traceroute de GravNet a Location server, y de Spycam a Shrouded Step, con cada router intermedio indicado.

\subsubsection*{GravNet $\rightarrow$ Location server}
\evid{Q2_conectividad/ping_gravnet_location.png}{Ping de GravNet a Location server.}
\evid{Q2_conectividad/tracert_gravnet_location.png}{Traza de GravNet a Location server.}

\begin{tabular}{@{}cll@{}}
  \toprule
  Salto & Router & IP observada \\
  \midrule
  1 & Deadlock & 172.16.24.1 \\
  2 & Cypher & 1.0.0.2 \\
  3 & Neon & 5.0.0.2 \\
  4 & Gekko & 8.0.0.2 \\
  5 & Reyna & 12.0.0.2 \\
  6 & Sage & 11.0.0.2 \\
  7 & Brimstone & 13.0.0.2 \\
  8 & Location server & 192.168.50.2 \\
  \bottomrule
\end{tabular}

\medskip
El paquete recorre 7 routers y atraviesa los cinco sistemas autónomos de la red (AS 10, 20, 30, 40 y 50) hasta llegar al Location server. El ping responde
sin pérdidas y el TTL alterna entre 120 y 121, lo que indica que los paquetes no siguen siempre el mismo camino: algunos atraviesan un router más que otros.
Esto es coherente con el reparto del tráfico entre dos rutas de igual métrica que Packet Tracer muestra en Deadlock y Cypher (ver pregunta 3). Entre Cypher y Neon la traza no muestra
routers intermedios, porque en cada salto Packet Tracer informa la IP de la interfaz por la que responde y no por la que entra el paquete (Ver preguntas 3 y 4)


\subsubsection*{Spycam $\rightarrow$ Shrouded Step}
\evid{Q2_conectividad/ping_spycam_shrouded.png}{Ping de Spycam a Shrouded Step.}
\evid{Q2_conectividad/tracert_spycam_shrouded.png}{Traza de Spycam a Shrouded Step.}

\begin{tabular}{@{}cll@{}}
  \toprule
  Salto & Router & IP observada \\
  \midrule
  1 & Cypher & 172.16.192.1 \\
  2 & Killjoy & 3.0.0.2 \\
  3 & Fade & 6.0.0.2 \\
  4 & Viper & 7.0.0.2 \\
  5 & Omen & 9.0.0.2 \\
  6 & Shrouded Step & 4.93.112.4 \\
  \bottomrule
\end{tabular}

\medskip
El paquete sale del AS 10 por Killjoy, cruza Fade (AS 20) y Viper (AS 30) y entra al AS 40 por Omen, para llegar hasta Shrouded Step en 6 saltos. Killjoy prefiere esta salida porque 
aprende la red de la LAN de Omen por eBGP desde Fade, con distancia administrativa 20, mientras que la misma red también le llega por OSPF con distancia 110. El router instala la de 
menor distancia (Ver la tabla de Killjoy en la pregunta 6).